\section{Phase-velocity slices}\label{sec:slices}

This section shows that each branch (\ref{eq:branch}) is a continuum of passive scalars, or slices, labelled by the parallel phase velocity. The Alfv\'enic fields $\Phi$ and $\Psi$ are given functions of $\boldsymbol r$ and $t$; (\ref{eq:rmhd}) is not used.

\subsection{Without the Landau coupling}\label{sl:nocoupling}

For $\Lambda=\infty$ the operator $L$ is multiplication by $u$, and (\ref{eq:branch}) holds separately at every $u$:
\begin{equation*}
  \partial_tG+\{\Phi,G\}+su\,\gradpar G=0 .
\end{equation*}
Since $\gradpar=\partial_z+\{\Psi,\cdot\}$ and $\{f,G\}=(\zhat\times\nabla_\perp f)\cdot\nabla G$, this is the advection equation $\partial_tG+\boldsymbol U_u\cdot\nabla G=0$ with the velocity
\begin{equation}\label{eq:Uu}
  \boldsymbol U_u=\zhat\times\nabla_\perp(\Phi+su\Psi)+su\,\zhat=\uE+su\,\bhat,\qquad\nabla\cdot\boldsymbol U_u=0 .
\end{equation}
At each $u$, $G$ is a passive scalar, carried by the $\boldsymbol E\times\boldsymbol B$ flow and streaming along the perturbed field at the speed $su$, which is the motion of the guiding centre of a free ion with $v_\parallel=u\vth$. Values of $G$ at different $u$ do not interact. This is the exact equation for the Laguerre coefficients $G_l$ with $l\ge2$ (S09, eqs~210 and~212).

\subsection{With the Landau coupling}\label{sl:transformsec}

For finite $\Lambda$, $LG=uG-(uF_M/\Lambda)\,M_0[G]$ ties every velocity to the density $M_0[G]$, through the parallel electric field and the mirror force, so $G$ at fixed $u$ is no longer a passive scalar. The slices survive in other variables. The reason is that $L$ acts on velocity only and is the same at every point and at every wavenumber, because the field relations (\ref{mod:fields}) are algebraic. A change of velocity variables that diagonalises $L$ therefore does so everywhere at once, and it passes through the advection and the streaming unchanged.

The argument is clearest for the hierarchy (\ref{mod:hierarchy}) truncated at $m=M$ with the zero closure, in which $L$ is the constant $(M+1)\times(M+1)$ matrix $A$. Let $\boldsymbol w_j$ be a left eigenvector of $A$, $\boldsymbol w_j^{\mathsf T}A=u_j\boldsymbol w_j^{\mathsf T}$. Since $\boldsymbol w_j$ does not depend on position or time, it passes through $\partial_t$, $\partial_z$ and both brackets, and multiplying (\ref{mod:hierarchy}) by $\boldsymbol w_j^{\mathsf T}$ gives for $a_j=\boldsymbol w_j^{\mathsf T}g$
\begin{equation*}
  \partial_ta_j+\{\Phi,a_j\}+su_j\,\gradpar a_j
  =\partial_ta_j+\{\Phi+su_j\Psi,\,a_j\}+su_j\,\partial_za_j=0 .
\end{equation*}
Each $a_j$ is a passive scalar, advected with the stream function $\Phi+su_j\Psi$ and streaming along $z$ at the speed $su_j$, as $G$ at fixed $u$ was for $\Lambda=\infty$, and nothing couples $a_j$ to the other $a_l$. The nonlinearity enters only through the bracket with $\Phi+su_j\Psi$, and the Landau coupling does not appear, because it has been absorbed into $\boldsymbol w_j$. For the $a_j$ to be a change of variables, the $u_j$ must be real and the $\boldsymbol w_j$ linearly independent, so that the $a_j$ determine $g$. Both hold because $\Pw A$ is symmetric (section~\ref{mod:hermite}): $A$ is self-adjoint in the free-energy product $\sum_mp_mg_mh_m$, which is positive definite when $p_0=1-1/\Lambda>0$, that is, for $\Lambda\notin[0,1]$. Section~\ref{sec:truncated} normalises the $a_j$ and treats both closures.

At $M=\infty$ the same steps need the eigenfunctions of $L$ itself. Self-adjointness carries over: $L$ is self-adjoint in the free-energy product
\begin{equation}\label{eq:fe}
  \fe{f}{g}=\int\frac{fg}{F_M}\,\dd u-\frac{M_0[f]\,M_0[g]}{\Lambda},
\end{equation}
which is $\sum_mp_mf_mg_m$ in the Hermite variables and is positive definite for $\Lambda\notin[0,1]$, as for every compressive branch of KRMHD ($\Lambda^+>1$, $\Lambda^-<0$; section~\ref{sec:model}). Unlike $A$, however, $L$ has no eigenfunctions of finite free energy: its spectrum is the whole real line and is continuous. Its generalised eigenfunctions are the Case--Van Kampen modes $F_u$ of the branch \citep{vanKampen1955,Case1959}, with $LF_u=uF_u$ for every real $u$. Each is a beam of particles at $v_\parallel=u\vth$, a $\delta$-function in velocity, dressed by the response of the Maxwellian, through the Landau coupling, to the density of the beam (appendix~\ref{app:spectral}). The projection $\Tr[G](u)=\fe{F_u}{G}$ on $F_u$ takes the place of $\boldsymbol w_j^{\mathsf T}g$, and its unit-weight form is the slice amplitude
\begin{equation}\label{eq:aunit}
  a_u=\sqrt{\varrho(u)}\;\fe{F_u}{G},\qquad \varrho(u)=\frac{F_M(u)}{|\diel(u+\ii0)|^2},
\end{equation}
where $\diel$ is the dielectric function of the branch, which contains the plasma dispersion function, and $\varrho$ is the spectral density of the branch, positive with $\int\varrho\,\dd u=1/p_0$. For $\Lambda=\infty$, $F_u$ is the bare beam, $\diel=1$, $\varrho=F_M$ and $a_u=G/\sqrt{F_M}$. The map $G\mapsto a_u$ is a fixed linear change of variables in velocity. It is invertible, divides the free energy among phase velocities without cross terms, and turns $L$ into multiplication by $u$:
\begin{equation}\label{sl:parseval}
  \fe{G}{G}=\int a_u^2\,\dd u,\qquad G\mapsto a_u\ \text{ takes }\ LG\mapsto u\,a_u
\end{equation}
(theorem~\ref{thm:spectrum}, which also gives $a_u$ in terms of $G$ and its Hilbert transform in $u$). Its inverse is, up to normalisation, the G-transform that \citet{MorrisonPfirsch1992} and \citet{Morrison2000} introduced for linearised Vlasov--Poisson; \citet{Heninger2018} applied it to electrostatic drift kinetics, where the transformed components remain coupled through the potential.

\subsection{The slice theorem}\label{sl:slicesec}

\begin{theorem}\label{thm:slice}
Assume that (i) $\Lambda\notin[0,1]$, the value $\Lambda=\infty$ included; (ii) $\Phi(\boldsymbol r,t)$ and $\Psi(\boldsymbol r,t)$ are real functions on $\mathbb T^3\times[0,t_1]$, with $\mathbb T^3$ a periodic box, whose gradients are continuous and Lipschitz in $\boldsymbol r$, uniformly in $t$; they need not satisfy (\ref{eq:rmhd}) and may obey resistive or forced RMHD, or be prescribed; (iii) $G$ and $uG$ are continuously differentiable functions of $(\boldsymbol r,t)$ with values in the space $\mathcal W$ of functions of $u$ with finite $\int G^2F_M^{-1}\,\dd u$ (appendix~\ref{app:spectral}). Then $G$ solves (\ref{eq:branch}) if and only if, at every $(\boldsymbol r,t)$ and for almost every $u$,
\begin{equation}\label{eq:slice}
  \partial_ta_u+\boldsymbol U_u\cdot\nabla a_u=\partial_ta_u+\{\Phi+su\Psi,\,a_u\}+su\,\partial_za_u=0,
\end{equation}
with $\boldsymbol U_u$ of (\ref{eq:Uu}).
\end{theorem}

\begin{proof}
This is the argument for the truncated hierarchy in section~\ref{sl:transformsec}, with the map $\mathcal A:G\mapsto a$ of (\ref{eq:aunit}) in place of $\boldsymbol w_j^{\mathsf T}$. The map $\mathcal A$ is bounded and injective (theorem~\ref{thm:spectrum}) and acts on $u$ only, through a kernel independent of $\boldsymbol r$ and $t$, so it commutes with $\partial_t$, with spatial derivatives and with multiplication by functions of $(\boldsymbol r,t)$, hence with $\{\Phi,\cdot\}$, $\{\Psi,\cdot\}$ and $\gradpar$. With $\mathcal ALG=u\,\mathcal AG$ this gives $\mathcal A[\partial_tG+\{\Phi,G\}+s\,\gradpar(LG)]=\partial_ta+\{\Phi,a\}+su\,\gradpar a$, and $u\,\gradpar a_u=u\,\partial_za_u+\{u\Psi,a_u\}$ gives (\ref{eq:slice}). Since $\mathcal A$ is linear and injective, $G$ solves (\ref{eq:branch}) if and only if $a$ solves (\ref{eq:slice}).
\end{proof}

With the Landau coupling, therefore, each slice is again a passive scalar carried by $\boldsymbol U_u$, in three dimensions, and slices with different $u$ do not interact. The parallel electric field and the mirror force, which enter through the Landau coupling, only dress the slices: the distribution at fixed $u$ is not a passive scalar, but its projection $a_u$ on $F_u$ is. Since $\boldsymbol U_u$ is divergence-free, each slice keeps its own free energy $\int a_u^2\,\dd^3\boldsymbol r$, and by (\ref{sl:parseval}) the nonlinearity moves no free energy between phase velocities (section~\ref{sec:invariants}). In terms of the Elsasser potentials, the stream function of the slice at $u$ is
\begin{equation}\label{eq:elsasserpot}
  \Phi+su\Psi=\tfrac12\big[(1+su)\,z^++(1-su)\,z^-\big],
\end{equation}
so the slices at $su=-1$ and $+1$ are advected by $z^-$ alone and by $z^+$ alone, and the slice at $u=0$ by $\Phi$ alone (section~\ref{sl:special}). The sectors have different transforms, but $\boldsymbol U_u$ depends only on $u$, so at each $u$ the slices of $G^+$, $G^-$ and every $G_l$ with $l\ge2$ are carried by the same flow:
\begin{equation}\label{eq:sameflow}
  \big(\partial_t+\boldsymbol U_u\cdot\nabla\big)\big(a^+_u,\,a^-_u,\,a^{(2)}_u,\,a^{(3)}_u,\dots\big)=0 .
\end{equation}
The compressive sector of KRMHD is therefore a continuum of passive scalars labelled by phase velocity. Section~\ref{sec:truncated} shows that a truncated Hermite hierarchy decomposes in the same way, and figure~\ref{fig:snapshots} (section~\ref{sec:numerics}) shows three slices of a GANDALF run that start from the same density pattern and are carried apart by their own flows.

\subsection{Special slices}\label{sl:special}

The slice $u=0$ is advected by $\uE$ alone. If $G$ is H\"older continuous in $u$ near $u=0$, (\ref{eq:transform}) gives $\Tr[G](0)=p_0\,G(0)/F_M(0)$, so this slice is the distribution at $v_\parallel=0$, which obeys $\partial_tG+\{\Phi,G\}=0$ there because $(LG)(0)=0$. This is the KRMHD form of CMM's eq.~(2.22), by which in gyrokinetics the distribution at $v_\parallel=0$ is advected by the gyroaveraged potential alone. Section~\ref{sec:invariants} relates the energy of this slice to the imaginary part of CMM's complex helicity, and section~\ref{sec:truncated} gives its truncated counterpart.

\begin{proposition}\label{prop:elsasser}
Under the hypotheses of theorem~\ref{thm:slice}, the slices with $su=-1$ and $su=+1$, that is, with $v_\parallel=-v_A$ and $+v_A$, obey
\begin{equation}\label{sl:elsasser}
  (\partial_t-\partial_z)\,a_u+\{z^-,a_u\}=0\quad(su=-1),\qquad
  (\partial_t+\partial_z)\,a_u+\{z^+,a_u\}=0\quad(su=+1).
\end{equation}
The slice $su=-1$ moves toward $-z$ with $z^+$ and is advected by $z^-$ alone, and the slice $su=+1$ moves toward $+z$ with $z^-$ and is advected by $z^+$ alone. In a truncated hierarchy the same holds for a slice of theorem~\ref{thm:trunc} whose speed $su_j$ is $-1$ or $+1$.
\end{proposition}

\begin{proof}
By (\ref{eq:elsasserpot}), $\Phi+su\Psi$ is $z^-$ at $su=-1$ and $z^+$ at $su=+1$. Substituting in (\ref{eq:slice}) gives (\ref{sl:elsasser}).
\end{proof}

The two slices exist at every $\beta_i$, at $u=\pm\beta_i^{-1/2}$ (in a truncated hierarchy only when $\beta_i=u_j^{-2}$ for an eigenvalue $u_j$), whereas the fluid slow waves interact exclusively with the counter-propagating Alfv\'en waves only for $\beta\gg1$ \citetext{S09, \S2.5, eq.~30 and footnote~15; \citealp[\S2.4]{LithwickGoldreich2001}}. The fluid slow waves and the entropy mode of S09 (their eqs~23 and~29--30) obey (\ref{eq:slice}) at $su=\mp c$, with $c$ the slow-wave speed in units of $v_A$, and at $u=0$; this is a formal correspondence, not a derived collisional limit of the slices.

The equations (\ref{eq:rmhd}) and (\ref{eq:branch}) are invariant under the reflection
\begin{equation}\label{eq:mirror}
  \Phi'(\boldsymbol r)=\Phi(\tilde{\boldsymbol r}),\qquad
  \Psi'(\boldsymbol r)=-\Psi(\tilde{\boldsymbol r}),\qquad
  G'(\boldsymbol r,u)=G(\tilde{\boldsymbol r},-u),\qquad
  \tilde{\boldsymbol r}=(x,y,-z),
\end{equation}
which reverses $z$ and $v_\parallel$, exchanges $z^+$ and $z^-$, and maps the slice at $u$ to the slice at $-u$, $a'_u(\boldsymbol r)=a_{-u}(\tilde{\boldsymbol r})$. The free energy $\Wc$ is therefore even under (\ref{eq:mirror}), and $\Gamma$ and $\Hph$, which weight the slice energies with $u$ and $1/u$ (section~\ref{sec:invariants}), are odd, so they vanish in every state that is its own mirror image.

\subsection{Dissipation, collisions and forcing}\label{sl:forcingsec}

Let the right-hand side of (\ref{eq:branch}) be $-\mathcal BG-\mathcal CG+F$, where the linear operators $\mathcal B$ and $\mathcal C$ act on $\boldsymbol r$ alone and on $u$ alone, respectively, and $F$ is a source. The proof of theorem~\ref{thm:slice} gives
\begin{equation}\label{eq:sliceforcing}
  \partial_ta_u+\boldsymbol U_u\cdot\nabla a_u=-\mathcal B\,a_u-\int\hat{\mathcal C}(u,u')\,a_{u'}\,\dd u'+\sqrt{\varrho(u)}\;\Tr[F](u),
\end{equation}
where $\hat{\mathcal C}$ is the kernel of $\mathcal A\,\mathcal C\,\mathcal A^{-1}$. Damping at a rate that depends only on $k_\perp$ or $k_z$, applied with the same coefficient to every Hermite moment, is of the form $\mathcal B$ and damps each slice separately. With $e_m=H_m(u)F_M(u)/\sqrt{2^mm!}$ and $r_m(u)$ the $m$th Hermite coefficient of $F_u$, a polynomial of degree $m$ with $r_0=1$ and $r_1=\sqrt2\,u$, one has $\Tr[e_m](u)=p_m\,r_m(u)$, and the functions $\phi_m=\sqrt{p_m\varrho}\;r_m$ form an orthonormal basis of $L^2(\mathbb R,\dd u)$. For $\mathcal Ce_m=\nu_me_m$, as for Lenard--Bernstein collisions ($\nu_m=\nu m$) or hypercollisions ($\nu_m=\nu(m/M)^n$), the kernel is $\hat{\mathcal C}(u,u')=\sum_m\nu_m\,\phi_m(u)\,\phi_m(u')$, which couples different slices unless all $\nu_m$ are equal; collisions are the only term in (\ref{eq:sliceforcing}) that exchanges free energy between slices. Forcing of $g_0$ by a field $f(\boldsymbol r,t)$, $F=f\,e_0$, gives every slice the same spatial pattern, with amplitude $(1-1/\Lambda)\sqrt{\varrho(u)}\,f$.

\begin{remark}\label{sl:remlit}
Parts of this structure are known in other settings. For $\Lambda=\infty$, theorem~\ref{thm:slice} restates S09's eqs~(210) and~(212), which describe the Laguerre components with $l\ge2$ as a passively mixed, undamped ballistic-type mode; the transform adds the two branches, absorbing their Landau coupling into the slices. \citet[section~IV]{Heninger2018} applied the G-transform to the electrostatic drift-kinetic ion equation with Boltzmann electrons in a slab \citep[see also][chapter~13]{Heninger2020}, where quasineutrality also makes the kernel independent of the wavenumber; their transform is ours with $\Lambda=-T_i/(ZT_e)$, the limit of $\Lambda^-$ as $\beta_i\to0$. Their transformed components stream along the uniform guide field and share one $\boldsymbol E\times\boldsymbol B$ flow but remain coupled through the potential, which is proportional to $\int g\,\dd v_\parallel$; their equation also conserves the energy of each component in three dimensions, although they do not state this. In linear problems at a single wavenumber the transform was used by \citet{Plunk2013}, for a forced Landau problem, and by \citet{PlunkParker2014}, who noted that without forcing any $u$-weighted integral of the squared transformed distribution is conserved, and \citet{Plunk2015} used a transform that commutes with $\boldsymbol E\times\boldsymbol B$ advection to obtain a quadratic invariant at each velocity, and from it the nonlinear stability of a quasi-two-dimensional drift-kinetic model without parallel streaming.

Invariants labelled by a velocity variable are known in two dimensions. In two-field models of collisionless reconnection, two Lagrangian invariants are advected by the generalised stream functions $\phi\pm(\rho_s/d_e)\psi$ \citep{Schep1994,Cafaro1998,Grasso2001}, and in their drift-kinetic version the distribution at each value of the parallel canonical momentum is a Lagrangian invariant \citep{Liseikina2004,Pegoraro2004}. Two-dimensional gyrokinetics has an invariant for every pair $(v_\perp,v_\parallel)$ \citetext{S09, appendix~F.1; \citealp[\S3]{Plunk2010}}, and \citet{ZhuHammett2010} used per-velocity quadratic invariants that survive Galerkin truncation to construct its absolute equilibria; the low-beta equations of \citet{ZoccoSchekochihin2011} have Lagrangian invariants parametrised by $v_\parallel$. In each of these models the terms with parallel gradients break the invariants in three dimensions.
\end{remark}
